%-------------------------
% Resume in Latex
% Author : Jake Gutierrez
% Based off of: https://github.com/sb2nov/resume
% License : MIT
%------------------------

\documentclass[letterpaper,11pt]{article}

\usepackage{latexsym}
\usepackage[empty]{fullpage}
\usepackage{titlesec}
\usepackage{marvosym}
\usepackage[usenames,dvipsnames]{color}
\usepackage{verbatim}
\usepackage{enumitem}
\usepackage[hidelinks]{hyperref}
\usepackage{fancyhdr}
\usepackage[english]{babel}
\usepackage{tabularx}
\input{glyphtounicode}

\usepackage{fontawesome} % Required for FontAwesome icons


%----------FONT OPTIONS----------
% sans-serif
% \usepackage[sfdefault]{FiraSans}
% \usepackage[sfdefault]{roboto}
% \usepackage[sfdefault]{noto-sans}
% \usepackage[default]{sourcesanspro}

% serif
% \usepackage{CormorantGaramond}
% \usepackage{charter}


\pagestyle{fancy}
\fancyhf{} % clear all header and footer fields
\fancyfoot{}
\renewcommand{\headrulewidth}{0pt}
\renewcommand{\footrulewidth}{0pt}

% Adjust margins
\addtolength{\oddsidemargin}{-0.5in}
\addtolength{\evensidemargin}{-0.5in}
\addtolength{\textwidth}{1in}
\addtolength{\topmargin}{-.5in}
\addtolength{\textheight}{1.0in}

\raggedbottom
\raggedright
\setlength{\tabcolsep}{0in}

% Sections formatting
\titleformat{\section}{
  \vspace{-4pt}\scshape\raggedright\large
}{}{0em}{}[\color{black}\titlerule \vspace{-5pt}]

% Ensure that generate pdf is machine readable/ATS parsable
\pdfgentounicode=1

%-------------------------
% Custom commands
\newcommand{\resumeItem}[1]{
  \item\small{
    {#1 \vspace{-2pt}}
  }
}

\newcommand{\resumeSubheading}[4]{
  \vspace{-2pt}\item
    \begin{tabular*}{0.97\textwidth}[t]{l@{\extracolsep{\fill}}r}
      \textbf{#1} & #2 \\
      \textit{\small#3} & \textit{\small #4} \\
    \end{tabular*}\vspace{-7pt}
}

\newcommand{\resumeSubSubheading}[2]{
    \item
    \begin{tabular*}{0.97\textwidth}{l@{\extracolsep{\fill}}r}
      \textit{\small#1} & \textit{\small #2} \\
    \end{tabular*}\vspace{-7pt}
}

\newcommand{\resumeProjectHeading}[2]{
    \item
    \begin{tabular*}{0.97\textwidth}{l@{\extracolsep{\fill}}r}
      #1 & #2 \\
    \end{tabular*}\vspace{-7pt}
}

\newcommand{\resumeExperienceheading}[2]{
  \vspace{-2pt}\item
    \begin{tabular*}{0.97\textwidth}[t]{l@{\extracolsep{\fill}}r}
      \textbf{#1} & #2 
    \end{tabular*}\vspace{-7pt}
}


\newcommand{\resumeSubItem}[1]{\resumeItem{#1}\vspace{-4pt}}

\renewcommand\labelitemii{$\vcenter{\hbox{\tiny$\bullet$}}$}

\newcommand{\resumeSubHeadingListStart}{\begin{itemize}[leftmargin=0.15in, label={}]}
\newcommand{\resumeSubHeadingListEnd}{\end{itemize}}
\newcommand{\resumeItemListStart}{\begin{itemize}}
\newcommand{\resumeItemListEnd}{\end{itemize}\vspace{-5pt}}

%-------------------------------------------
%%%%%%  RESUME STARTS HERE  %%%%%%%%%%%%%%%%%%%%%%%%%%%%


\begin{document}

\begin{center}
    \textbf{\Huge \scshape Juan Ignacio Biancuzzo} \\ \vspace{1pt} 
    { \Large Software Engineer $\sim$ Estudiante de ingeniería } \\

    { \small
        \faGlobe \href{https://juanbiancuzzo.github.io/JuanBiancuzzo/}{ portafolio} $|$
        \faEnvelope \href{mailto:biancuzzojuanignacio@gmail.com}{ biancuzzojuanignacio@gmail} $|$
        \faGithub \href{https://github.com/JuanBiancuzzo}{ github.com/JuanBiancuzzo} $|$
        \faLinkedinSquare \href{https://www.linkedin.com/in/juan-ignacio-biancuzzo-367611231/}{/juan-ignacio-biancuzzo}
    }
    
\end{center}

\section{Resumen profesional}
    \resumeSubHeadingListStart
        \vspace{-2pt}\item Soy un estudiante de Ingeniería en Informática, altamente motivado y comprometido con el aprendizaje continuo. Poseo una sólida experiencia en la realización de proyectos y me destaco en habilidades clave como la comunicación efectiva, la resolución de problemas y la empatía, lo que me permite asumir nuevos desafíos con gran disposición y eficacia
    \resumeSubHeadingListEnd

%-----------PROGRAMMING SKILLS-----------
\section{Habilidades técnicas}
 \begin{itemize}[leftmargin=0.15in, label={}]
    \small{\item{
     \textbf{Lenguajes de programación}{: C, C++, C\#, Rust, Java, Python y SQL} \\
     \textbf{Herramientas}{: Git, Docker, Unity, Pandas, Numpy y Matplotlib} \\
     \textbf{Lenguajes}{: \textbf{Ingles} - C1, \textbf{ Español} - Nativo} \\
    }}
 \end{itemize}

 %-----------EDUCATION-----------
\section{Educación}
  \resumeSubHeadingListStart
    \resumeSubheading
      {Universidad de Buenos Aires (UBA)}{}
      {Estudios universitarios en las carreras de ingeniería en informática y electrónica}{Marzo 2019 -- Actualidad}
    \resumeSubheading
      {Colegio de LaSalle}{}
      {Estudios primarios y secundarios}{Julio. 2002 -- Noviembre 2018}
  \resumeSubHeadingListEnd

%-----------PROJECTS-----------
\section{Proyectos}
    \resumeSubHeadingListStart
        \resumeProjectHeading {\textbf{Solo} $|$ \emph{Python y Godot}}{\href{https://github.com/ptourne/uncert-titan}{github.com/uncert-solo}}
            \resumeItemListStart
                \resumeItem{\textbf{Participación destacada en la Hackathon Internacional Space Apps Challenge de la NASA}: Competí en un evento global de dos días enfocado en soluciones para la exploración espacial}
                \resumeItem{\textbf{Desarrollo y lanzamiento de \href{https://sologame.ar}{sologame.ar}}: Contribuí a la conceptualización y creación de un juego interactivo diseñado para educar sobre desafíos y soluciones en la exploración espacial}
                \resumeItem{\textbf{Desempeño como Developer}: Colaboré estrechamente con un compañero para implementar las funcionalidades del juego, optimizando la experiencia del usuario y el rendimiento técnico}
                \resumeItem{\textbf{Liderazgo en la organización del equipo}: Coordiné las actividades del grupo, asigné tareas y gestioné el tiempo de manera eficiente para cumplir con los plazos estrictos del hackathon}
                \resumeItem{\textbf{Gestión efectiva del tiempo y recursos}: Supervisioné el progreso del proyecto, asegurando la entrega exitosa y a tiempo del juego, a pesar de las limitaciones del evento}
            \resumeItemListEnd

      \resumeProjectHeading {\textbf{Marching cubes - Unity package} $|$ \emph{C\# y HLSL}}{\href{https://github.com/It-is-not-only-me/com.itisnotonlyme.marching-cubes}{github.com/marching-cubes}}
            \resumeItemListStart
                \resumeItem{\textbf{Visualización 3D eficiente y detallada}: Conseguí una visualización eficiente y detallada al implementar el algoritmo de Marching Cubes utilizando Compute Shaders en Unity, lo que resultó en un rendimiento gráfico optimizado y una representación precisa de datos volumétricos}
                \resumeItem{\textbf{Desafíos técnicos complejos}: Superé desafíos complejos y amplié mis habilidades al enfrentar y resolver problemas avanzados fuera de mi zona de confort, lo que resultó en un producto final robusto y de alta calidad}
                \resumeItem{\textbf{Rendimiento gráfico y la calidad visual}: Mejoré significativamente el rendimiento gráfico al desarrollar y optimizar Fragment y Vertex Shaders, lo que resultó en una experiencia visual más fluida y atractiva en la aplicación}
            \resumeItemListEnd
    \resumeSubHeadingListEnd

%-----------EXPERIENCE-----------
\section{Experiencia}
  \resumeSubHeadingListStart

    \resumeExperienceheading
      {Ayudante de la materia de Algoritmos y Programación I}{Marzo 2022 -- Diciembre 2023}
      \resumeItemListStart
        \resumeItem{\textbf{Asistencia y tutoría}: Guié a estudiantes en la comprensión y aplicación de conceptos clave de algoritmos y programación, ayudándolos a mejorar sus habilidades técnicas y a progresar en su formación profesional}
        \resumeItem{\textbf{Soporte académico}: Brindé apoyo individualizado a estudiantes, respondiendo a preguntas y aclarando dudas para facilitar su aprendizaje y mejorar su desempeño en la materia}
      \resumeItemListEnd

    \resumeExperienceheading
      {Ayudante de la materia de Taller de Programación I}{Agosto 2023 -- Actualidad}
      \resumeItemListStart
        \resumeItem{\textbf{Desarrollo de trabajos prácticos}: Participé en la creación y diseño de trabajos prácticos, asegurando que estos reflejen y promuevan las cualidades esenciales de colaboración en equipo entre los estudiantes}
        \resumeItem{\textbf{Soporte en la enseñanza}: Asistí en la implementación de las prácticas, proporcionando orientación y apoyo a los estudiantes para garantizar la comprensión de los conceptos y el éxito en las tareas asignadas}
    \resumeItemListEnd
      
% -----------Multiple Positions Heading-----------
%    \resumeSubSubheading
%     {Software Engineer I}{Oct 2014 - Sep 2016}
%     \resumeItemListStart
%        \resumeItem{Apache Beam}
%          {Apache Beam is a unified model for defining both batch and streaming data-parallel processing pipelines}
%     \resumeItemListEnd
%    \resumeSubHeadingListEnd
%-------------------------------------------

  \resumeSubHeadingListEnd

%-------------------------------------------
\end{document}
